% BHU PYQ -- Transcribed & Corrected 100% Accurate
% Paper: CSCSE21: Mobile Application Development
% Exam: B.Sc. (Hons.) Semester II (NEP) Examination 2024-25
% Category: SEC (Department of Computer Science)
% Ref Code: CECONF/N/2024-25/21

\documentclass[11pt,a4paper]{article}
\usepackage[a4paper,top=2.5cm,bottom=2.5cm,left=2.5cm,right=2.5cm,headheight=14pt]{geometry}
\usepackage{amsmath,amssymb}
\usepackage[T1]{fontenc}
\usepackage[utf8]{inputenc}
\usepackage{lmodern,microtype,fancyhdr,enumitem,hyperref,lastpage,parskip,listings}

\hypersetup{colorlinks=true,urlcolor=black,linkcolor=black,pdftitle={CSCSE21: Mobile Application Development -- BHU Sem II 2024-25 (NEP SEC)}}
\pagestyle{fancy}\fancyhf{}
\fancyhead[L]{\small   \textit{Banaras Hindu University}}
\fancyhead[R]{\small   \textit{CSCSE21: Mobile App Development (End Term SEC)}}
\fancyfoot[C]{\small Page  \thepage\ of \pageref{LastPage}}


\newlist{parts}{enumerate}{2}
\setlist[parts,1]{label=(\alph*),leftmargin=*,itemsep=4pt,topsep=2pt}
\setlist[parts,2]{label=(\roman*),leftmargin=*,itemsep=2pt,topsep=2pt}

\newcommand{\pts}[1]{\hfill   \textit{[#1]}}

\lstset{
  language=SQL,
  basicstyle=\small tfamily,
  keywordstyle=\bfseries,
  showstringspaces=false,
  frame=single,
  tabsize=2,
  breaklines=true
}


\begin{document}


\begin{center}
  {\large\bfseries BANARAS HINDU UNIVERSITY}\\[2pt]
  {
\normalsize B.Sc. (Hons.) Semester II (NEP) Examination 2024-25}\\[6pt]
  \rule{0.8\linewidth}{0.5pt}\\[6pt]
  {\large\bfseries DEPARTMENT OF COMPUTER SCIENCE}\\[4pt]
  {\large\bfseries Skill Enhancement Course (SEC)}\\[4pt]
  {\large\bfseries Paper Code: CSCSE21 --- Mobile Application Development}\\[4pt]
  {
\normalsize Time: 2:30 Hours\hfill Maximum Marks: 70}\\[2pt]
  {\small Paper Code Ref: CECONF/N/2024-25/21}
\end{center}

\rule{\linewidth}{0.4pt}\smallskip



\noindent   \textit{Instruction: Write your Roll No. at the top immediately on the receipt of this question paper.}\\[2pt]



\noindent   \textit{Note: Question 1 is compulsory. Attempt any five from the remaining six questions.}
\bigskip




\noindent   \textbf{Question 1}   \textit{Choose the correct option from the following questions:}\pts{$10  \times 2 = 20$}

\begin{parts}
  \item What is Android?
    
\begin{parts}
      \item A mobile device
      \item An operating system
      \item A web browser
      \item A programming language
    \end{parts}

  \item Which file is the entry point for an Android application?
    
\begin{parts}
      \item   \texttt{MainActivity.kt}
      \item   \texttt{AndroidManifest.xml}
      \item   \texttt{strings.xml}
      \item   \texttt{build.gradle}
    \end{parts}

  \item What does APK stand for?
    
\begin{parts}
      \item Android Phone Kit
      \item Android Package Kit
      \item Android Primary Key
      \item Application Package Key
    \end{parts}

  \item Which language is officially supported for Android development?
    
\begin{parts}
      \item Python
      \item Swift
      \item Java
      \item C++
    \end{parts}

  \item Which language has been recommended by Google for Android development since 2019?
    
\begin{parts}
      \item Kotlin
      \item C\#
      \item JavaScript
      \item Ruby
    \end{parts}

  \item What is the purpose of   \texttt{R.java} file?
    
\begin{parts}
      \item To store app settings
      \item To manage database
      \item To link XML with Java/Kotlin code
      \item To debug code
    \end{parts}

  \item Which component is used to display a UI screen in Android?
    
\begin{parts}
      \item Broadcast Receiver
      \item Service
      \item Activity
      \item Fragment
    \end{parts}

  \item Which method is called when an activity is first created?
    
\begin{parts}
      \item   \texttt{onStart()}
      \item   \texttt{onCreate()}
      \item   \texttt{onResume()}
      \item   \texttt{onRestart()}
    \end{parts}

  \item Which layout arranges elements in a single column or row?
    
\begin{parts}
      \item ConstraintLayout
      \item GridLayout
      \item FrameLayout
      \item LinearLayout
    \end{parts}

  \item Which of the following is used for local data storage in Android?
    
\begin{parts}
      \item Firebase
      \item MySQL
      \item SQLite
      \item PostgreSQL
    \end{parts}
\end{parts}

\medskip\rule{\linewidth}{0.2pt}\medskip




\noindent   \textbf{Question 2}   \textit{Define any five of the following terms:}\pts{$5  \times 2 = 10$}

\begin{parts}
  \item APK
  \item Gradle in Android Studio
  \item Layout in Android Studio
  \item   \texttt{AndroidManifest.xml}
  \item Activity in Android
  \item Intent in Android development
\end{parts}

\medskip\rule{\linewidth}{0.2pt}\medskip




\noindent   \textbf{Question 3}   \textit{Write short answer type questions:}\pts{$2  \times 5 = 10$}

\begin{parts}
  \item Describe the components of Android application development.
  \item Explain the use of SQLite in Android with an example.
\end{parts}

\medskip\rule{\linewidth}{0.2pt}\medskip




\noindent   \textbf{Question 4}\pts{10} \\
Explain any five UI components of Android application development.

\medskip\rule{\linewidth}{0.2pt}\medskip




\noindent   \textbf{Question 5}\pts{$2  \times 5 = 10$}

\begin{parts}
  \item What is a content provider in Android? What are the operations of the content provider?
  \item Explain the working of the content provider.
\end{parts}

\medskip\rule{\linewidth}{0.2pt}\medskip




\noindent   \textbf{Question 6}   \textit{Write the SQLite query for the following statements:}\pts{$5  \times 2 = 10$}

\begin{parts}
  \item Create a   \texttt{student} table with attributes:   \texttt{ID},   \texttt{name},   \texttt{age},   \texttt{gender},   \texttt{department}, and   \texttt{marks}.
  \item Write an SQLite query to insert a new record with details: Name: Raj, Age: 21, Gender: male, Department: Computer Science, Marks: 85.
  \item Write a query to display names of the students from   \texttt{student} table.
  \item Write a query to display all student records from table.
  \item Write a query to update the marks of the student with   \texttt{ID = 3} to   \texttt{90}.
\end{parts}

\medskip\rule{\linewidth}{0.2pt}\medskip




\noindent   \textbf{Question 7}\pts{$2  \times 5 = 10$}

\begin{parts}
  \item Write a Kotlin program with a function   \texttt{add} that returns the sum of two numbers.
  \item Write a Kotlin program to print numbers from 1 to 10 using a loop statement.
\end{parts}

\vfill\rule{\linewidth}{0.4pt}

\begin{center}\small
\href{https://bhu-exam.vercel.app/}{Click here to visit our website}\\[4pt]
\href{https://me-aryan.vercel.app/}{Click here to visit our contributor}
\end{center}

\end{document}
